\section{Lecture 22: Toward Berkovich Spaces --- Localization over
Topological Spaces and the Stack of Norms}
\label{lec:22}

This is the penultimate lecture of the course. Clausen returns to the
programme announced in \cref{lec:1}: to see that the traditional frameworks of
analytic geometry sit inside the theory of analytic stacks. Having already
treated adic spaces (\cref{lec:10}) and objects like the Tate curve
(\cref{lec:19,lec:21}), he now moves toward \emph{Berkovich spaces}. The
lecture has two halves. First a piece of setup: several inequivalent notions of
how to \emph{localize} an analytic stack over a topological space, phrased
through locales and the six-functor formalism, and how the notion used earlier
(for metrizable finite-dimensional compact Hausdorff spaces) sits among them.
Second, the Berkovich spectrum $\MBerk(A)$ of a Banach ring, and its
reincarnation as an analytic stack $\mathcal N$ of \emph{norms}; the bulk of the
time is spent computing the ``image'' of $\mathcal N$ over $\mathbb Z$, which
turns out to be an \emph{extended} Berkovich space of $\mathbb Z$, with the
archimedean branch prolonged past the usual absolute value to a compactifying
point.

Throughout, ``ring'' means commutative ring and everything is (light)
condensed.

\begin{remark}[Derived Huber pairs]
\label{rmk:22:derived-huber}
Asked about the sheafiness defect of adic spaces (\cref{lec:10}) and whether one
must pass to some notion of \emph{derived} Huber pair / derived adic space,
Clausen said this is still in progress and possibly admits several
formulations; a definition of derived Huber pairs appears in Rodr\'iguez
Camargo's recent paper on the analytic de Rham stack in rigid geometry
\cite{camargo2024analyticrhamstackrigid}.
\end{remark}

\subsection{Localizing an analytic stack over a topological space}
\label{subsec:22:localization}

We have already seen (\cref{thm:20:compact-hausdorff}) one way to localize an
analytic stack over a space. Let $S$ be a metrizable, finite-dimensional
compact Hausdorff space --- or a locally compact Hausdorff space locally of
this form --- viewed (\cref{ex:19:condensed-anima}) as an analytic stack. Given
$X\in\AnStk$ one may simply ask for a map
\begin{equation}
\label{eq:22:map-to-S}
f\colon X\longrightarrow S \qquad\text{in }\AnStk .
\end{equation}
As discussed, this equips $D(X)$ with a family of idempotent algebras indexed by
the closed subsets $Z\subseteq S$ (the $i_{Z*}$ of constant sheaves) --- or,
dually, idempotent coalgebras $j_!\mathbf 1_U$ indexed by the open complements
--- and these algebra objects localize $D(X)$ over $S$. For the comparison with
Berkovich spaces the metrizability/finite-dimensionality hypotheses, though
satisfied in practice, are an annoyance, so we introduce weaker notions of
localization along a topological space and relate them to
\eqref{eq:22:map-to-S}.

To any analytic stack $X$ one attaches its underlying \emph{locale}
\begin{equation*}
\Loc(X)\ :=\ \{\,Y\subseteq X\,\}\ =\ \{\text{monomorphisms into }X\text{ in
}\AnStk\},
\end{equation*}
the poset of subobjects of $X$. This is the general construction of the
underlying locale of an ($\infty$-)topos, retaining only the monomorphisms; the
topos axioms guarantee the requisite unions and finite intersections (the
category of analytic stacks has all limits and colimits, so one forms unions as
images of coproducts). A monomorphism means, concretely, that the diagonal
$Y\to Y\times_X Y$ is an isomorphism.

\begin{remark}[A possible size issue that does not matter]
\label{rmk:22:size}
It is not clear (Clausen did not investigate) whether $\Loc(X)$ is a set;
$\AnStk$ is not presentable, only closed under the same exactness properties as
an $\infty$-topos. This is harmless for what follows: in the definitions below
one maps \emph{into} $\Loc(X)$ from an honest locale $S$, and, $S$ and $X$ being
fixed, the resulting data is a set --- there are no nontrivial automorphisms of
any one of these maps, so no anima appear.
\end{remark}

Now fix a locale $S$ (for a topological space, its poset of opens; we abuse
$U\subseteq S$ to mean $U$ an element of this poset). There are three
increasingly strong notions of localizing $X$ over $S$.
\begin{enumerate}
\item A map of locales (geometric morphism)
\begin{equation}
\label{eq:22:map-of-locales}
\pi\colon \Loc(X)\longrightarrow S ,
\end{equation}
i.e.\ for every open $U\subseteq S$ a monomorphism $\pi^{-1}U\subseteq X$,
sending unions to unions and finite intersections to finite intersections.
\item The stronger requirement that each inclusion of opens $U\subseteq V
\subseteq S$ be sent to an \emph{open immersion}
\begin{equation*}
\pi^{-1}U\ \subseteq\ \pi^{-1}V\ \subseteq\ X
\end{equation*}
from the perspective of the (extended) six-functor formalism of \cref{lec:21}
--- that is, a shriekable, cohomologically smooth monomorphism. (For a
monomorphism the relative dualizing object is canonically the unit, so
cohomological smoothness is exactly the condition of being an open immersion.)
Equivalently, one factors \eqref{eq:22:map-of-locales} through the quotient
locale $\Loc^{\varphi}(X)$ of those monomorphisms which are open immersions,
\begin{equation}
\label{eq:22:factor-through-phi}
\Loc(X)\longrightarrow \Loc^{\varphi}(X)\longrightarrow S .
\end{equation}
\item If $S$ is metrizable, finite-dimensional and compact Hausdorff, one may
ask for an actual map $f\colon X\to S$ in $\AnStk$ as in
\eqref{eq:22:map-to-S}.
\end{enumerate}

\begin{proposition}
\label{prop:22:implications}
$(3)\Rightarrow(2)\Rightarrow(1)$.
\end{proposition}

\begin{proof}[Indication]
$(2)\Rightarrow(1)$ is tautological. For $(3)\Rightarrow(2)$ one checks on the
analytic stack attached to $S$ that every inclusion of opens is an open
immersion in the six-functor sense: this is because our six-functor formalism,
restricted to (locally) compact Hausdorff spaces, recovers the usual one
(\cref{lec:21}), and the property of being an open immersion is stable under
base change. (Upgrading $(1)$ to $(2)$ is just the check that certain inclusions
are open immersions; upgrading $(2)$ to the strongest notion $(3)$, on the other
hand, imposes a genuine \emph{connectivity} condition on the idempotent algebras,
as before.)
\end{proof}

\begin{remark}[Unions of open immersions]
\label{rmk:22:union-open}
For $\Loc^{\varphi}(X)$ in \eqref{eq:22:factor-through-phi} to be a locale one
needs a union of open immersions to be an open immersion. This holds: a union
along open immersions is a cover by shriekable maps, so it acquires a
lower-shriek functor, and cohomological smoothness may be checked locally. It
follows from the extension procedure for six-functor formalisms of
\cref{lec:21}.
\end{remark}

\begin{example}[Huber pairs and the adic spectrum]
\label{ex:22:huber-locale}
Let $(R,R^+)$ be a Huber pair, with associated solid analytic ring
$(R,R^+)^\solid$ (\cref{def:9:solid-structure}), and set $X=\AnSpec((R,R^+)^
\solid)$. Then $\Loc(X)$ localizes over Huber's topological space
$\Spa(R,R^+)$ of continuous valuations: for a rational open
$U=U(f_1,\dots,f_n/g)$ the preimage is again an affine analytic stack, the one
enforcing $g$ invertible and $f_1/g,\dots,f_n/g$ power-bounded (solid),
\begin{equation*}
\pi^{-1}U\ =\ \AnSpec\bigl(R\langle f_1,\dots,f_n/g\rangle^\solid\bigr),
\end{equation*}
which is a monomorphism (on derived categories, a localization). This refines
\cref{thm:9:solid-sheaf}: the proof given there in fact showed the stronger
statement that any cover of a rational open in $\Spa(R,R^+)$ pulls back under
$\pi$ to an open cover in the six-functor sense.
\end{example}

\begin{remark}[Rational opens need not be open immersions]
\label{rmk:22:warning-rational-open}
In general it is \emph{not} true that a rational open pulls back to an open
immersion. Consider
\begin{equation*}
(\mathbb Z_p,\mathbb Z_p)^\solid\longrightarrow(\mathbb Q_p,\mathbb Z_p)^\solid,
\end{equation*}
i.e.\ inverting $p$ while keeping the ring of integral elements. On derived
categories,
\begin{equation*}
D\bigl((\mathbb Z_p,\mathbb Z_p)^\solid\bigr)\longrightarrow
D\bigl((\mathbb Q_p,\mathbb Z_p)^\solid\bigr)
\end{equation*}
is the inclusion of the full subcategory on which $p$ acts invertibly (the map
$p\colon X\xrightarrow{\sim}X$). From the six-functor perspective this is a
\emph{closed} immersion --- a proper monomorphism --- not an open one: only the
underlying ring changes, the analytic ring structure is unchanged.
\end{remark}

\begin{remark}[The Tate case, and a modified topology]
\label{rmk:22:tate-fix}
If $R$ is Tate, this pathology does not arise: inverting a function is a rising
union of rational subsets cut out by inequalities,
\begin{equation*}
\{f\neq0\}\ =\ \bigcup_{\varepsilon}\ \{\,|f|\ge\varepsilon\,\},
\end{equation*}
each of which \emph{is} open. (This subset is not quasicompact, though it is a
quasicompact rational domain in the adic space --- the discrepancy is that the
map to $\Spa(A,A^+)$ is not literally the naive one, being built via
$\Spa$ of the rational localizations and Huber's retraction.) There is, in
addition, another topology on $\Spa(R,R^+)$ with the same constructible subsets
--- incomparable to Huber's, a spectral space, the opposite of one of Huber's
topologies --- in which Zariski-open subsets are re-declared closed, so that
open pulls back to open even outside the Tate case. This modification is
discussed in the ``GAGA redux'' chapter of Clausen--Scholze's
\emph{Lectures on Analytic Geometry} \cite{clausen2020analytic}.
\end{remark}

\begin{remark}[Passing to inverse limits]
\label{rmk:22:inverse-limits}
If $(S_i)_{i\in I}$ is an inverse system of compact Hausdorff spaces, then to
give compatible maps of locales $\Loc(X)\to S_i$ for all $i$ is the same as to
give a single map
\begin{equation*}
\Loc(X)\longrightarrow \varprojlim_i S_i ,
\end{equation*}
since inverse limits of compact Hausdorff spaces agree with the corresponding
limits of locales here; the same holds for $\Loc^{\varphi}$. Combined with the
fact that \emph{every} compact Hausdorff space is an inverse limit of
metrizable finite-dimensional ones --- potentially in many different ways, but
often with a natural choice, as will occur for Berkovich spaces --- this lets
one bootstrap from the metrizable finite-dimensional case (where $(3)$ is
available) to general compact Hausdorff spaces, at the cost of
finite-dimensionality. Asking for the strongest notion, an actual map
$X\to\varprojlim_i S_i$ of analytic stacks, still requires the connectivity
condition ($\Loc^{\varphi}$-locally all idempotent algebras connective) and is
not automatic; Clausen expects it to be genuinely stronger, though produced no
counterexample.
\end{remark}

\subsection{Berkovich spectra: recollection}
\label{subsec:22:berkovich-recollection}

We recall the classical definition to fix the target.

\begin{definition}[Banach ring]
\label{def:22:banach-ring}
A \emph{Banach ring} $(A,|\cdot|)$ is a commutative ring $A$ together with a map
$|\cdot|\colon A\to\mathbb R_{\ge0}$ satisfying
\begin{equation*}
\text{(1) } |0|=0, \qquad
\text{(2) } |x+y|\le|x|+|y|, \qquad
\text{(3) } |xy|\le|x|\,|y|, \qquad
\text{(4) } A=0 \text{ or } |1|=1,
\end{equation*}
together with $|{-1}|=|1|$ (equivalently the triangle inequality also for
$x-y$), and such that $A$ is complete with respect to the metric
$d(x,y)=|x-y|$ (completeness in particular forces $|f|=0\iff f=0$).
\end{definition}

\begin{example}
\label{ex:22:banach-examples}
\leavevmode
\begin{enumerate}
\item For $S$ a compact Hausdorff space, $A=C(S;\mathbb C)$ with the sup-norm
\begin{equation*}
|f|\ =\ \sup_{s\in S}|f(s)|_{\mathrm{usual}} .
\end{equation*}
\item For $R$ a complete Tate--Huber ring with pseudo-uniformizer (topologically
nilpotent unit) $\pi\in R$, a norm $|\cdot|\colon R\to\mathbb R_{\ge0}$ with
$|\pi|=\tfrac12$, with $|R_0|\le1$ for a ring of definition $R_0$, and
multiplicative in $\pi$:
\begin{equation*}
|\pi f|\ =\ |\pi|\,|f| \qquad(f\in R).
\end{equation*}
(It is multiplicative when one factor is $\pi$, not in general.) E.g.\ $\mathbb
Q_p$ with $\pi=p$, or $\overline{\mathbb Q}_p$ with $\pi=p$.
\end{enumerate}
\end{example}

\begin{definition}[Berkovich spectrum]
\label{def:22:berkovich-spectrum}
To a Banach ring $(A,|\cdot|)$ Berkovich attaches
\begin{equation}
\label{eq:22:berkovich}
\MBerk(A,|\cdot|)\ =\ \bigl\{\,|\cdot|_x\colon A\to\mathbb R_{\ge0}\ \text{a
multiplicative seminorm with } |f|_x\le|f|\ \forall f\in A\,\bigr\},
\end{equation}
the set of multiplicative seminorms bounded by the given norm; ``multiplicative''
includes $|1|_x=1$ (so the point is nonzero) and $|0|_x=0$. It is topologized as
a closed subset of a product of intervals,
\begin{equation*}
\MBerk(A,|\cdot|)\ \hookrightarrow\ \prod_{f\in A}[0,|f|],\qquad
|\cdot|_x\ \longmapsto\ \bigl(|f|_x\bigr)_{f\in A},
\end{equation*}
and is compact Hausdorff. For $A=0$ it is empty.
\end{definition}

\begin{example}
\label{ex:22:berkovich-CS}
For $A=C(S,\mathbb C)$ the seminorms $|f|_x=|f(x)|$, $x\in S$, are points of
$\MBerk(A)$; if one moreover demands that $|\cdot|_x$ restrict to the usual norm
on $\mathbb C$, these are \emph{all} the points (Gelfand's theorem), recovering
$S$.
\end{example}

\subsection{The analytic stack of norms}
\label{subsec:22:norms}

The goal is to make the ``same definition'' as \eqref{eq:22:berkovich} --- the
set of multiplicative seminorms bounded by $|\cdot|$ --- but inside analytic
stacks rather than topological spaces, using the notion of a \emph{norm} on an
analytic ring from \cref{lec:20} (\cref{def:20:norm}).

\begin{definition}[Norm; recollection]
\label{def:22:norm}
For $X\in\AnStk$, a \emph{norm} on $X$ is a map of analytic stacks
\begin{equation*}
N\colon\ \mathbb P^1_X\ \longrightarrow\ [0,\infty]
\end{equation*}
which is multiplicative on $N^{-1}\bigl((0,\infty)\bigr)$, satisfies
$N(T^{-1})=N(T)^{-1}$ (compatibility of inversion on $\mathbb P^1$ with
$\lambda\mapsto\lambda^{-1}$ on $[0,\infty]$), and satisfies the condition on
how the basic disk $P$ sits: $P$ lies between the open unit disk
$N^{-1}\bigl([0,1)\bigr)$ and the closed unit disk $N^{-1}\bigl([0,1]\bigr)$
(\cref{def:20:norm}, Axiom~4). Here $T$ is the coordinate on $\mathbb P^1$.
\end{definition}

\begin{notation}[The stack of norms]
\label{not:22:stack-of-norms}
Write $\mathcal N(X)$ for the set of norms on $X$. This is genuinely a set, as it
is a set of maps in the locally small category $\AnStk$. Norms satisfy descent,
so $X\mapsto\mathcal N(X)$ is represented by an analytic stack, again denoted
$\mathcal N$: thus $\mathcal N(X)=\Homop_{\AnStk}(X,\mathcal N)$.
\end{notation}

\begin{remark}[The gaseous base is a cover]
\label{rmk:22:gaseous-cover}
There is a cover, in the Grothendieck topology on analytic stacks,
\begin{equation}
\label{eq:22:norms-cover}
\AnSpec\bigl(\mathbb Z[\widehat q^{\pm1}]_\gasx\bigr)\ \twoheadrightarrow\
\mathcal N ,
\end{equation}
classifying the universal norm on an analytic ring $(R,D(R))$ equipped with an
element $q\in\ur R$ whose norm is exactly $\tfrac12$: a priori $N(q)$ is a map
$\Spec(R)\to[0,\infty]$ with image possibly an interval, and one demands that it
factor through the singleton $\{\tfrac12\}$, a stringent condition whose
universal solution is the gaseous base ring $\mathbb Z[\widehat q^{\pm1}]_\gasx$
of \cref{lec:14,lec:20}. The map \eqref{eq:22:norms-cover} is a cover because,
given any norm, one may locally find such an element $q$
(\cref{lem:20:small-element}).
\end{remark}

\subsection{The image of \texorpdfstring{$\mathcal N$}{N}: an extended
Berkovich space of \texorpdfstring{$\mathbb Z$}{Z}}
\label{subsec:22:image}

We now explore what $\mathcal N$ looks like. Over an arbitrary analytic ring the
only ring elements available are the integers, so for $n\in\mathbb Z$ we get
$N(n)\colon\Spec(R)\to[0,\infty]$, and hence a map of locales
\begin{equation*}
\Loc(\mathcal N)\ \xrightarrow{\ \pi\ }\ \prod_{n\in\mathbb Z}[0,\infty] .
\end{equation*}
Thus $\mathcal N$ localizes over this (large, non-finite-dimensional) product.

\begin{definition}[Image]
\label{def:22:image}
The \emph{image} of $\mathcal N$ is the complement of the largest open subset
$U\subseteq\prod_{n\in\mathbb Z}[0,\infty]$ with $\pi^{-1}U=\emptyset$.
\end{definition}

We will see that the fibre over any point of the image is non-empty, so at least
in this case the image is closed. Recall first the classical picture.

\begin{remark}[The Berkovich space of $\mathbb Z$]
\label{rmk:22:MZ}
The Berkovich spectrum $\MBerk(\mathbb Z)$ (all points avoiding $\infty$)
embeds in $\prod_{n\in\mathbb Z}[0,\infty)$ as the multiplicative norms
$|\cdot|\colon\mathbb Z\to[0,\infty)$ with $|1|=1$, $|0|=0$ and $|x+y|\le
|x|+|y|$. As a topological space it is a tree
(\cref{fig:22:MZ}): a central point, the trivial norm $|\cdot|_0$ (value $1$ on
everything nonzero); an archimedean branch ending at the usual absolute value
$|\cdot|_\infty$, with $|\cdot|_\infty^{1/2}$ at its midpoint; and, for each
prime $p$, a branch of the powers $|\cdot|_p^{\alpha}$, $\alpha\in(0,\infty)$,
whose far end ($\alpha\to\infty$) is the pullback of the trivial norm on
$\mathbb F_p$ (value $0$ on multiples of $p$, $1$ elsewhere), the usual
$p$-adic norm ($|p|=1/p$) sitting partway along. It is a small,
one-dimensional subspace of the huge product.
\end{remark}

% Provenance: board 29 (00:11:07) --- Berkovich tree of Z, redrawn cleanly.
\begin{figure}[ht]
\centering
\begin{tikzpicture}[scale=1.0,>=stealth]
  \coordinate (O) at (0,0);
  % archimedean branch to the right
  \coordinate (Ah) at (2.4,-0.3);   % midpoint |.|^{1/2}
  \coordinate (Ae) at (4.8,-0.6);   % end |.|_infty
  \draw (O) -- (Ae);
  \fill (Ah) circle (1.4pt);
  \fill (Ae) circle (1.6pt);
  % p-branch up-right
  \coordinate (Pp) at (1.3,1.5);    % |.|_p
  \coordinate (Pe) at (1.9,2.3);    % end F_p
  \draw (O) -- (Pe);
  \fill (Pp) circle (1.4pt);
  \fill (Pe) circle (1.6pt);
  % l-branch up-left (another prime)
  \coordinate (Le) at (-1.6,2.0);
  \draw (O) -- (Le);
  \fill (Le) circle (1.6pt);
  % more branches (schematic)
  \draw (O) -- (-2.3,0.9);  \fill (-2.3,0.9) circle (1.6pt);
  \node at (-1.2,-1.0) {$\cdots$};
  \fill (O) circle (2.0pt);
  \node[below left] at (O) {$|\cdot|_0$};
  \node[below] at (Ah) {$|\cdot|_\infty^{1/2}$};
  \node[right] at (Ae) {$|\cdot|_\infty$};
  \node[right] at (Pp) {$|\cdot|_p$};
  \node[above] at (Pe) {$\mathbb F_p$};
\end{tikzpicture}
\caption{The Berkovich tree $\MBerk(\mathbb Z)$: trivial norm at the centre, one
branch per prime ending at $\mathbb F_p$, and an archimedean branch ending at
the usual absolute value.}
\label{fig:22:MZ}
\end{figure}

\begin{theorem}[The image is an extended Berkovich space]
\label{thm:22:image}
The image of $\mathcal N$ over $\mathbb Z$ is a larger subset than
$\MBerk(\mathbb Z)$: it agrees with the tree on every non-archimedean branch,
but the archimedean branch is prolonged past $|\cdot|_\infty$ all the way to
$\infty$ and then one-point compactified. In symbols
\begin{equation}
\label{eq:22:image-formula}
\mathrm{Im}(\mathcal N)\ =\ \Bigl(\MBerk(\mathbb Z)^{\mathbb R_{\ge1}}\times
\mathbb R_{>0}\Bigr)^{+},
\end{equation}
the one-point compactification of $\MBerk(\mathbb Z)$ with the exponent-rescaling
action of $\mathbb R_{\ge1}$ extended to $\mathbb R_{>0}$; the compactifying
``strange point'' is the norm with
\begin{equation*}
N(0)=0,\quad N(\pm1)=1,\quad N(n)=\infty\ \text{ for } n\neq0,\pm1 .
\end{equation*}
\end{theorem}

\begin{remark}
\label{rmk:22:triangle-fails}
In particular the triangle inequality can fail on $\mathrm{Im}(\mathcal N)$: at
the strange point $N(1)=1$ but $N(2)=\infty$; and already $|\cdot|_\infty^\alpha$
for $\alpha>1$, a genuine new point of the prolonged archimedean branch,
violates it. In \eqref{eq:22:image-formula} the superscript $\mathbb R_{\ge1}$ is
Clausen's shorthand for the exponent-rescaling action $N\mapsto N^{\alpha}$ by
$\alpha\ge1$ applied to $\MBerk(\mathbb Z)$: this action does nothing on the
non-archimedean branches (each is already a full ray) but extends the archimedean
branch beyond $|\cdot|_\infty$, after which the $(-)^{+}$ one-point
compactification adds the single point at infinity.
\end{remark}

To justify \cref{thm:22:image} we analyse the branches one at a time via the
function $N(p)$ for a fixed prime $p$.

\begin{proposition}[Interior of the $p$-adic branch]
\label{prop:22:p-branch}
For the map $N(p)\colon\mathcal N\to[0,\infty]$,
\begin{equation*}
\mathcal N_{0<|p|<1}\ :=\ N(p)^{-1}\bigl((0,1)\bigr)\ =\ \Spec\bigl(\mathbb
Q_{p}^{\gasx}\bigr)\times(0,1).
\end{equation*}
For a fixed $\lambda\in(0,1)$ the fibre $\Spec(\mathbb Q_p^{\gasx})\times
\{\lambda\}$ carries a normed analytic ring structure whose notion of
overconvergent functions on the disk of radius $r$ is the usual one from
non-archimedean geometry, for the normalization of the $p$-adic absolute value
with $|p|=\lambda$.
\end{proposition}

\begin{proof}[Indication]
The universal analytic stack equipped with a variable of norm in $(0,1)$ is
$\Spec(\mathbb Z[\widehat q^{\pm1}]_\gasx)\times(0,1)$ (\cref{rmk:22:gaseous-cover}).
Imposing that the variable become $p$ --- modding out by $q-p$ --- yields, as in
\cref{lec:14}, the gaseous analytic ring structure on $\mathbb Q_p$; the second
factor $(0,1)$ is unchanged.
\end{proof}

\begin{proposition}[Interior of the archimedean branch]
\label{prop:22:arch-branch}
The locus where $|p|$ is strictly between $1$ and $\infty$ is
\begin{equation*}
\mathcal N_{1<|p|<\infty}\ =\ \Spec\bigl(\mathbb R^{\gasx}\bigr)\times(0,1),
\end{equation*}
and it is independent of $p$. For fixed $\lambda\in(0,1)$ the universal norm on
$\Spec(\mathbb R^{\gasx})\times\{\lambda\}$ is given by the usual overconvergent
holomorphic functions of complex geometry, with respect to the power
$|\cdot|_\infty^{\alpha}$ of the usual absolute value where $\alpha$ is chosen so
that
\begin{equation*}
|\tfrac1p|_\infty^{\alpha}\ =\ \lambda .
\end{equation*}
\end{proposition}

\begin{proof}[Indication]
That $1<|p|<\infty$ means $p$ is invertible and $0<|1/p|<1$; applying the
argument of \cref{prop:22:p-branch} to the variable $1/p$ gives
$\Spec(\mathbb Z[\widehat q^{\pm1}]_\gasx)\times(0,1)$, and setting $q=1/p$
produces the gaseous analytic ring structure on $\mathbb R$ (\cref{lec:14}).
Independence of $p$ is visible from the description, which never referred to $p$;
this is the analytic-stacks incarnation of Ostrowski's classification --- once
$|p|$ is pinned in this range the norm of every other integer is determined.
\end{proof}

\begin{remark}
\label{rmk:22:overconvergence-washes-out}
The subtlety of exactly which summability condition defines $\mathbb R^{\gasx}$
(some quasi-exponential decay, \cref{rmk:14:sums}) --- a space sitting strictly
between the overconvergent functions on the closed unit disk and the holomorphic
functions on the open disk --- washes out once one passes to \emph{over}\-convergent
functions on a disk of a given radius: the answer is just the usual ring of
overconvergent holomorphic functions, up to a scaling of the absolute value
coming from the chosen value of $N(q)$.
\end{remark}

\subsubsection*{Triangle inequalities}

The image computation rests on the following universal claims, in which
$2$ may be replaced by any integer $>1$.

\begin{proposition}[Universal triangle inequalities]
\label{prop:22:triangle}
Let $(R,N)$ be a normed analytic ring.
\begin{enumerate}
\item If $N(2)\le1$, then $N$ satisfies the non-archimedean triangle inequality:
\begin{equation}
\label{eq:22:triangle-na}
\mathcal N_{|2|\le1}\ \subseteq\ \MBerk{}^{\mathrm{na}}(\mathbb Z).
\end{equation}
\item If $N(2)\le2$, then $N$ satisfies the usual triangle inequality.
\item If $N(2)<\infty$, then there is a constant $C>0$ with
\begin{equation}
\label{eq:22:triangle-quasi}
|x+y|\ \le\ C\,\bigl(|x|+|y|\bigr).
\end{equation}
\end{enumerate}
\end{proposition}

These are statements about maps of analytic stacks, not about literal
real-valued norm functions. For instance \eqref{eq:22:triangle-quasi} is the
assertion that for $A,B<\infty$ the addition map, which carries the product of
disks $\mathbb P^1_{|T|\le A}\times\mathbb P^1_{|S|\le B}$ into the affine line
$\mathbb A^1_R\subseteq\mathbb P^1_R$ (\cref{lem:20:mult-well-defined}), fits
into a commuting square

% Provenance: composed from Clausen's verbal description, board 38 (01:31:39).
% The addition map (T,S) |-> T+S is the vertical arrow into P^1_R (it lands in
% A^1_R); the top arrow is the induced factorization of N(T+S) through the
% bounded interval.
\begin{equation*}
\begin{tikzcd}[column sep=large, row sep=large]
\mathbb P^1_{|T|\le A}\times\mathbb P^1_{|S|\le B}
  \arrow[r] \arrow[d, "{(T,S)\mapsto T+S}"']
& {[0,\,C(A+B)]} \arrow[d, hook] \\
\mathbb P^1_R \arrow[r, "N"'] & {[0,\infty] .}
\end{tikzcd}
\end{equation*}

\begin{remark}[Proof of the triangle claims]
\label{rmk:22:triangle-proof}
By the cover \eqref{eq:22:norms-cover} one may assume there is $q$ with
$0<N(q)<1$, so that one works over the gaseous base $\mathbb Z[\widehat q^{\pm1}]
_\gasx$. In the universal case, the condition $N(2)\le1$ makes one live over an
idempotent algebra in $D(\mathbb Z[\widehat q^{\pm1}]_\gasx)$, namely the
overconvergent closed unit disk modded by $T=2$,
\begin{equation*}
\mathcal O_{\mathbb Z[\widehat q^{\pm1}]_\gasx}\bigl(\{|T|\le1\}\bigr)\big/(T-2).
\end{equation*}
A calculation identifies this, on underlying rings, with
\begin{equation*}
\mathbb Z(\!(q)\!)\ \cap\ \mathcal O^{\mathrm{hol}}_{q,0}
\end{equation*}
--- Laurent series with integer coefficients that are meromorphic at the origin,
i.e.\ converge on some unspecified open disk around $0$. This description is
manifestly independent of $2$, and one checks the non-archimedean triangle
inequality by hand (it amounts to the existence of a map of idempotent algebras
lying over the addition map on the polynomial ring). The claim $N(2)\le2$ is
proved similarly, setting $T^{-1}=\tfrac12$ instead: one gets a version of
functions converging on the open disk in which the Laurent tails force
convergence out to the boundary; it in turn implies the $N(2)<\infty$ claim.
\end{remark}

\begin{remark}[Proof of \cref{thm:22:image}]
\label{rmk:22:image-proof}
The image is computed by stratifying along $N(2)\colon\mathcal N\to[0,\infty]$
into $\{|2|\le1\}$, $\{1<|2|<\infty\}$ and $\{|2|=\infty\}$: a closed stratum
and its open complement need not form a cover, but they still detect
isomorphisms --- in particular emptiness --- so it suffices to work over each
stratum and take the union of the images. Over $\{|2|\le1\}$,
\eqref{eq:22:triangle-na} places us in the non-archimedean Berkovich space; over
$\{1<|2|<\infty\}$ we are in the (already-classified) archimedean locus
$\Spec(\mathbb R^\gasx)\times(0,1)$; over $\{|2|=\infty\}$ one needs $N(n)=\infty$
for all $n\neq0,\pm1$, which follows because $2$ was arbitrary (if some other
$n$ were finite, one would land in a lower stratum). The reverse containment ---
that each point is realized by an actual normed analytic ring --- holds because
the interiors were realized above, the central trivial norm by non-archimedean
geometry over a Laurent series ring $\mathbb Z(\!(q)\!)$ with $q$ of trivial
norm, and the end-points by non-archimedean geometry over $\mathbb F_p$.
\end{remark}

\subsection{The stack of norms as a stack; rescaling}
\label{subsec:22:closing}

\begin{remark}[Only Tate rings have norms; rescaling]
\label{rmk:22:tate-norms}
The assignment $R\mapsto\{\text{norms on }R\}$ is itself an analytic stack
$\mathcal N$, with the cover \eqref{eq:22:norms-cover}. Only Tate (analytic)
rings admit nontrivial norms; a discrete ring has no nontrivial theory of norms.
For a Tate--Huber ring a norm amounts to a choice of $|q|\in(0,1)$ for a
pseudo-uniformizer, and the space of norms is a torsor under the rescaling
$N\mapsto N^{\alpha}$ of a fixed one --- matching the rescaling on the Berkovich
spectrum, though for a global object it is a whole compatible family rather than
a single seminorm. Over $\Zsolid$ the reals vanish, so no norm with
$N(q)\in(0,1)$ exists there and $\mathcal N$ is a genuinely different kind of
base. Modding out by these exponential rescalings gives a quotient stack
$\mathcal N/\mathbb R_{>0}$, to which any Tate adic space admits a unique map.
This is to be taken up next time.
\end{remark}

\begin{remark}[The fibres are genuinely stacky at the ends]
\label{rmk:22:fibres-stacky}
On the interior of a branch the fibre feels like a space --- the unit interval
base-changed to a gaseous $\mathbb Q_p$ or $\mathbb R$
(\cref{prop:22:p-branch,prop:22:arch-branch}). But toward the outer end-points
(and the central point) it is genuinely a stack: over the boundary point of a
$p$-branch the fibre carries e.g.\ $\mathbb F_p(\!(q)\!)$ with its notion of
overconvergent functions on a disk, and rescaling the norm changes \emph{which}
disk of a given radius one is labelling --- so the fibre is not affine, and the
$\mathbb R_{>0}$-action, though split as a set, carries real geometric content.
\end{remark}

\begin{remark}[Administrative]
\label{rmk:22:admin}
There is no lecture on the Friday; Scholze and others give related talks in
person. The Berkovich discussion continues the following Wednesday, and the
last class is the Friday after --- the course concludes then (the IH\'ES
schedule listing two further weeks is to be corrected).
\end{remark}
